\setlength{\textfloatsep}{20.0pt plus 2.0pt minus 4.0pt}
\setlength{\intextsep}{12.0pt plus 2.0pt minus 2.0pt}
\section{Preliminaries}\label{sec:prelim}

This section sets up the shared definitions and external facts needed to state and prove the bridge, permutation, and Fourier-transform bounds of Sections~\ref{sec:bridge}--\ref{sec:main}. Section~\ref{sec:notation} collects global notation in Table~\ref{tab:notation}, and Section~\ref{sec:prelim-model} specifies the multitape Turing machine model, tape array layout, and fixed-point rounding rules. Section~\ref{sec:prelim-dft} fixes our normalized discrete Fourier transform conventions and the two-call real convolution identity, while Section~\ref{sec:prelim-oai} records the two black-box facts borrowed from~\cite{OAI26}. A first-time reader may skim Table~\ref{tab:notation} and Section~\ref{sec:prelim-oai} and return to them when cited.

\subsection{Notation}\label{sec:notation}

We begin by fixing basic conventions for logarithms, index ranges, and $\ell^\infty$ norms used across all sections. Logarithms are natural, written $\ln$, unless written $\log_2$. Following~\cite{OAI26}, we write $\lg n=\max\{\lceil\log_2 n\rceil,1\}$, so $\lg R=\log_2 R$ whenever $R\ge 2$ is a power of two. For a positive integer $N$ we write $[N]=\{0,\dots,N-1\}$. For a complex vector~$x$, $\|x\|_\infty=\max_j|x_j|$ denotes the largest complex modulus, and for a linear map~$A$, $\|A\|_\infty$ is the induced $\ell^\infty$ operator norm, equal to the maximum absolute row sum. For an array of polynomial records in $\Rr_r=\C[y]/(y^r+1)$, $\|\cdot\|_\infty$ denotes the largest modulus of any complex coefficient across the array. Table~\ref{tab:notation} summarizes global symbols.

\begin{table}[!t]
\centering
\small
\renewcommand{\arraystretch}{0.92}
\caption{Reference summary of global symbols and parameters used throughout the paper. Bridge-specific quantities are defined in Section~\ref{sec:bridge}, and finite-network parameters in Section~\ref{sec:prelim-oai}.}
\label{tab:notation}
\begin{tabular}{@{}>{\raggedright}p{0.27\textwidth}>{\raggedright\arraybackslash}p{0.67\textwidth}@{}}
\toprule
Symbol & Meaning\\
\midrule
$n,\ \kappa$ & bit length of each input integer; log-saving exponent in $O(n(\lg n)^{1-\kappa})$\\
$R=t_1\cdots t_d,\ k=\lg R$ & bridge transform length ($t_i\ge 2$ powers of two, $d$ dimension, $r=t_d=2^\ell$) and its binary logarithm\\
$M=R/r,\ \mathsf M(N),\ \mathsf m(N)$ & number of records; cost of the independent $N$-bit integer multiplier; $\mathsf m(N)=\mathsf M(N)/N$\\
$p,\ P$ & target precision of a bridge call (errors in units of $2^{-p}$); internal working precision of the bridge\\
$q,\ v=2^{-q}$ & dyadic grid precision of Sections~\ref{sec:walsh}--\ref{sec:direct-fft}, instantiated by the bridge at $q=P+2$; unit of that grid\\
$w$ & two's-complement component width; $w=q+2=P+4$ in the bridge schedule~\eqref{eq:schedule}\\
$\tw,\ \alpha$ & alias width parameter and Gaussian width parameter of the bridge kernel\\
$a\in\{2,\dots,d\},\ \nu\in\Z$ & spread-axis index of the bridge; alias index in Poisson summation\\
$g_a,\ G$ & negative half-bin carrier on axis $a$; Fourier transform of $\sinc(s)\,e^{-\pi s^2/\alpha^2}$\\
$f(\mathbf u),\ A_a$ & scalar frequency index with digits $u_1,\dots,u_d$; $A_a=\prod_{i<a}t_i$, the cycle length of $B_a$\\
$c_0,\ c_1,\ c_{\max}$ & retained alias weights $G((u+\tfrac12)/t)$ and $G((u+\tfrac12)/t-1)$; $\max(c_0,c_1)$\\
$\cA_a,\ B_a$ & full alias operator on axis $a$; two-term cyclic operator (its two-alias truncation)\\
$\Rr_r=\C[y]/(y^r+1)$ & record ring specifying the $r$-coefficient record layout\\
$\Cg_q,\ \mathbb{D}_q$ & fixed-point grid $2^{-q}(\Z+i\Z)$; disk grid $\{z\in\Cg_q:|z|\le 1\}$\\
$Q_q,\ \Quant_q$ & componentwise truncation toward zero to $\Cg_q$; componentwise nearest rounding to $\Cg_q$\\
$V$ & logical bit volume of a stored array\\
$K=\lfloor d^c\rfloor,\ D$ & chunk width and number of active $K$-bit chunks ($1\le D\le d$) in a simultaneous butterfly layer\\
$A,\ q_{\mathrm{rad}},\ \Gamma$ & layer and stream interface symbols: prefix range (written $P$ in~\cite[Prop.~14]{OAI26}), fixed radix (written $q$ in~\cite{OAI26}), and self-delimiting header code~\cite[(1)]{OAI26}\\
$H_0,\ S$ & normalized butterfly $(u,v)\mapsto((u+v)/2,(u-v)/2)$ and phase gate $\diag(1,i)$\\
$h=100,\ m=h^3=10^6$ & ground-set size parameters of the fixed triple networks of~\cite{OAI26}\\
$W_b,s_b;\ W_c,s_c$ & role counts and recursive call counts of the bit and complex networks ($s_b<W_bm$, $s_c<W_cm$)\\
$\eta_b,\ \eta_c;\ \tau,\ \sigma$ & relative rank deficits $(Wm-s)/(Wm)$ and network exponents ($s_b/W_b<m^\tau$, $s_c/W_c<m^\sigma$)\\
$\beta,\ \nu_{\mathrm{net}},\ C_1$ & leaf-stopping exponent, guard degree $\min\{j\ge 1:m^j\ge B_{\mathrm{net}}\}=11$, and guard exponent of the layer\\
$\eps,\ c,\ \lambda,\ \lambda'$ & schedule and recurrence exponents; $d=\lfloor k^\eps\rfloor$, $K=\lfloor d^c\rfloor$\\
$S,\ L,\ B,\ K$ & also reused locally (defined at use): suffix range or pre-scale/series width~$S$, layout map or cycle/input length~$L$, bit-reversal or suffix range~$B$, general chunk width~$K$ (and $K_0,K_*$)\\
\bottomrule
\end{tabular}
\end{table}

\subsection{Machine model and fixed-point arithmetic}\label{sec:prelim-model}\label{sec:prelim-grids}

This subsection defines the multitape Turing machine cost model, the sequential tape representation of multidimensional arrays, and the fixed-point complex grids on which the algorithm operates. Our complexity bounds refer to the standard deterministic multitape Turing machine model of~\cite{SS71,HvdH21,OAI26} with one fixed finite alphabet and a fixed finite number of one-dimensional tapes. In each step the machine reads and writes at most one cell per tape and moves each head by at most one cell. The running time is the total number of steps, so every scan, head movement, rewind, counter update, and workspace cleanup is charged. The input is $x\#y$ with $x,y\in\{0,1\}^n$, and the output is the $2n$-bit binary representation of their product~\cite[\S1]{OAI26}.

Following~\cite[\S2]{OAI26}, an array on $[a_1]\times\cdots\times[a_t]$ with $a_i\ge 1$ is stored in lexicographic order with the first coordinate most significant, so addresses are implicit in tape position rather than stored per entry. Adjacent spectator fields, meaning address fields that the current operation leaves untouched, may be collapsed into one field whose range is their product. An address permutation~$f$ moves the entry at address~$x$ to $f(x)$, giving $(fA)(f(x))=A(x)$. The \emph{logical bit volume}~$V$ of an array is its total number of stored data bits; temporary copies on work tapes are charged in running time without replacing~$V$ in recursive bounds. Integer parameters in procedure headers use the self-delimiting code $\Gamma(v)=0^{b_v-1}\mathrm{bin}_{b_v}(v)$ of length $2b_v-1$, where $b_v=\lfloor\log_2 v\rfloor+1$ for $v\ge 1$~\cite[(1)]{OAI26}.

\paragraph{Independent integer multiplier.}
Ordinary products of $N$-bit integers are delegated to an independent deterministic multiplier of worst-case cost $\mathsf M(N)=N\,\mathsf m(N)$, such as Sch\"onhage--Strassen multiplication~\cite{SS71} with $\mathsf m(N)=O(\lg N\lg\lg N)$. It runs as a black-box subroutine on dedicated work tapes that are clean on return, leaving the caller's other tapes untouched and making no recursive call to Algorithm~1.

\paragraph{Grids, rounding, and faithful tables.}
All stored complex numbers lie on dyadic fixed-point grids. The next definition fixes the grid $\Cg_P$ and disk grid $\mathbb{D}_P$ at an arbitrary integer precision $P\ge 0$, and specifies how records of $r$ complex numbers in $\mathbb{D}_q$ are stored in two's-complement binary at precision $q\ge 0$. It also introduces the two rounding operators, truncation toward zero $Q_P$ and nearest rounding $\Quant_P$, whose bounds~\eqref{eq:rounding} govern where each is used below.

\begin{definition}[Fixed-point grids and rounding]\label{def:grid}
For an integer precision $P\ge 0$, define the \emph{fixed-point grid} $\Cg_P=\{2^{-P}(a+ib):a,b\in\Z\}$ (also called
the \emph{$P$-grid}) and the \emph{disk grid} $\mathbb{D}_P=\{z\in\Cg_P:|z|\le 1\}$. At precision $q\ge 0$, a
\emph{disk-grid record} is an element $\sum_{j<r}z_j y^j\in\Rr_r=\C[y]/(y^r+1)$ with every $z_j\in\mathbb{D}_q$, stored
in coefficient order $j=0,\dots,r-1$; a \emph{disk array} is an array of records with $\|z\|_\infty\le 1$. Each real or
imaginary component is stored as $2^{-q}$ times a signed $w$-bit two's-complement integer, most significant bit first,
with the real component before the imaginary one~\cite[Prop.~14]{OAI26}; any width $w\ge q+2$ represents every
disk-grid component, including $\pm 1$. For $z\in\C$ and any integer $P\ge 0$:
\begin{itemize}[leftmargin=*]
  \item $Q_P(z)$ truncates the real and imaginary parts of $z$ toward zero to $2^{-P}\Z$~\cite[(26)]{OAI26};
  \item $\Quant_P(z)$ rounds the real and imaginary parts of $z$ to a nearest element of $2^{-P}\Z$ (for concreteness
    with ties to even, though every bound we prove holds for any nearest-rounding tie rule).
\end{itemize}
Both operators act coefficientwise on records and arrays.
\end{definition}

For every $z\in\C$ and integer $P\ge 0$,
\begin{equation}\label{eq:rounding}
  |Q_P(z)-z|<\sqrt2\cdot 2^{-P},\qquad |Q_P(z)|\le |z|,\qquad |\Quant_P(z)-z|\le\tfrac{\sqrt2}{2}\cdot 2^{-P}.
\end{equation}
In particular, $Q_q$ maps the closed unit disk into $\mathbb{D}_q$, so any exact disk-preserving operator followed by $Q_q$ keeps every stored coefficient inside $\mathbb{D}_q$, whereas nearest rounding can step outside the unit disk near $\pm 1$. By contrast, inside the bridge's spreading and deconvolution stages, where norms are tracked explicitly, $\Quant_q$ halves the per-step rounding error. Precomputed transcendental entries for filter taps, solver parameters, and unit phases are never assumed to be exact nearest roundings. Instead, they are \emph{faithful} approximations computed by truncated rational series with certified error at most $2^{-q}$ for real entries and $2\cdot 2^{-q}$ for complex taps and disk-valued unit phases.

\subsection{Discrete Fourier transforms}\label{sec:prelim-dft}\label{sec:prelim-oai-transform}

This subsection fixes the normalization and sign convention of our one-dimensional and multidimensional Fourier transforms and records the identity that recovers a real cyclic convolution from two forward transforms. Dividing by the transform length~$R$ in the definition below ensures that every exact transform is an $\ell^\infty$ contraction, keeping intermediate entries bounded without rescaling across dimensions.

\begin{definition}[Normalized transforms]\label{def:dft}
For an integer $R\ge 1$, the normalized positive one-dimensional DFT $F_R\colon\C^R\to\C^R$ is
$(F_R x)_u=\frac1R\sum_{j\in[R]}x_j e^{2\pi iju/R}$. For $\mathbf t=(t_1,\dots,t_d)$ with $R=t_1\cdots t_d$, the
normalized positive tensor DFT $F_{\mathbf t}=F^+_{\mathbf t}=F_{t_1,\dots,t_d}$ acts on arrays indexed by
$[t_1]\times\cdots\times[t_d]$ via
\[
  (F_{\mathbf t}U)_{u_1,\dots,u_d}=\frac1R\sum_{j_1,\dots,j_d}U_{j_1,\dots,j_d}\,
  \exp\Bigl(2\pi i\sum_{i=1}^d\frac{j_i u_i}{t_i}\Bigr).
\]
\end{definition}

By the triangle inequality, $\|F_R\|_\infty\le 1$ and $\|F_{\mathbf t}\|_\infty\le 1$. The cyclic convolution $(a\circledast b)_j=\sum_{l\in[R]}a_l b_{(j-l)\bmod R}$ satisfies $F_R(a\circledast b)=R\,(F_R a)\odot(F_R b)$, where $\odot$ denotes entrywise multiplication. For real vectors $a,b\in\R^R$ and $z=a+ib$, the \emph{two-call identity} (cf.~\cite{MS}) recovers $a\circledast b$ from two forward transforms of length~$R$:
\begin{equation}\label{eq:two-call}
  a\circledast b=-\frac{R^2}{2}\,\mathrm{Im}\Bigl(F_R\bigl(\overline{(F_R z)\odot(F_R z)}\bigr)\Bigr).
\end{equation}
Indeed, $\sum_{u\in[R]}e^{2\pi i(j-k)u/R}$ equals $R$ for $j=k$ and $0$ otherwise, so $F_R(\overline{F_R w})=R^{-1}\overline w$ for every $w\in\C^R$. Applying this relation to $w=z\circledast z=(a\circledast a-b\circledast b)+2i(a\circledast b)$, for which $(F_R z)\odot(F_R z)=R^{-1}F_R w$, gives $F_R(\overline{(F_R z)\odot(F_R z)})=R^{-2}\overline w$, whose imaginary part is $-2R^{-2}(a\circledast b)$.

For completeness, and only for the Bluestein comparison in Appendix~\ref{app:bluestein}, we also record the synthetic transforms of~\cite[\S6]{OAI26}; a first-time reader may skip this paragraph. Over the negacyclic ring $\Rr_r=\C[y]/(y^r+1)$ with $r=2^\ell\ge 2$, multiplication by~$y$ is an $\ell^\infty$ isometry. For any power of two $t\mid 2r$, setting $\omega_t=y^{2r/t}\in\Rr_r$ defines contractions $F_t^\mp\colon\Rr_r^t\to\Rr_r^t$ by $(F_t^\mp z)_j=t^{-1}\sum_{k<t}\omega_t^{\mp jk}z_k$ satisfying $F_t^+ F_t^-=t^{-1}I$. With $B_t$ denoting $(\lg t)$-bit binary reversal on $[t]$, in-place decimation in frequency and its reverse compute $B_t F_t^-$ and $F_t^+ B_t^{-1}$. For $\mathbf t'=(t_1,\dots,t_{d-1})$ with $t_i\mid 2r$, we write $F_{\mathbf t'}^\pm=\bigotimes_{i=1}^{d-1}F_{t_i}^\pm$ and $B=\bigotimes_{i=1}^{d-1}B_{t_i}$. Our direct complex tensor FFT in Section~\ref{sec:direct-fft} uses no ring products or synthetic transforms in $\Rr_r$, keeping $\Rr_r$ only to specify the contiguous $r$-coefficient record format.

\paragraph{Motivation from two-dimensional fractional shifts.}
In dimension two ($R=t_1 r$), Section~\ref{sec:overview-bridge} read the Cooley--Tukey twiddle factor as a fractional shift of row~$j_1$ by $j_1/t_1$~\cite{MS} and spread each row by a sinc-Gaussian filter. Applying Poisson summation then relates $F_{t_1,r}$ of the spread array, on the same $t_1\times r$ grid, to a cyclic alias mixture of $F_R x$ along each column in which two aliases dominate. With a real kernel the two retained alias weights coincide at the Nyquist column $u_2=r/2$; \cite{MS} handles this by a separate alternating-sum identity, and the coincidence prevents cycle-length-independent conditioning in higher dimensions. Section~\ref{sec:bridge} removes both obstacles for general $d\ge 2$ via a negative half-bin carrier and descending axis spreading.

\subsection{Ingredients from \texorpdfstring{\cite{OAI26}}{[OAI26]}}\label{sec:prelim-oai}

This subsection states the two results from~\cite{OAI26} that our written proofs use as black boxes, and summarizes the general-parameter tape procedures derived from them in Appendix~\ref{app:parameter-transfer}. Appendix~\ref{app:oai} restates the borrowed statements in full and records the symbol-renaming map from~\cite{OAI26}: its prefix range~$P$ becomes~$A$, its fixed radix~$q$ becomes $q_{\mathrm{rad}}$, and its box size~$T$ corresponds to~$R$.

First, Fact~\ref{fact:oai-streams} packages five linear-time sequential tape-stream primitives on arrays of bit volume~$V$. These primitives cover digit permutations, stream splits and merges, cyclic row distributions into $W$ role streams, prime-radix modular shifts, and exact fixed-width binary arithmetic.

\begin{fact}[Elementary tape operations~{\cite[Lemmas~3--4]{OAI26}}]\label{fact:oai-streams}
Fix an integer radix $q_{\mathrm{rad}}\ge 2$. On a stored array of logical bit volume $V$ on a fixed number of
finite-alphabet tapes, the following operations run in $O(V)$ steps (including counter initialization, head resets, and
clearing visited workspace):
\begin{enumerate}[label=(\arabic*),nosep,leftmargin=*]
  \item applying any fixed permutation to the pair of digit values at two specified radix-$q_{\mathrm{rad}}$ address
    positions, or moving one specified radix-$q_{\mathrm{rad}}$ digit position to another while preserving the order of
    intervening positions;
  \item splitting an array on a specified radix-$q_{\mathrm{rad}}$ digit, applying a fixed pointwise map to the aligned
    streams, and merging them back;
  \item within each prefix of an array on $[A]\times[R_{\mathrm{row}}]\times[B]$, distributing whole rows cyclically
    into a fixed number $W$ of role streams and, when $W\mid R_{\mathrm{row}}$, merging them in any fixed order of row
    classes;
  \item when $q_{\mathrm{rad}}$ is prime and a target field $y\in[q_{\mathrm{rad}}^b]$ follows control fields of
    $O(b)$ digits each across arbitrary spectators, shifting $y\mapsto y+a(\mathrm{controls})\pmod{q_{\mathrm{rad}}^b}$
    for any offset computable in $b^{O(1)}$ time (in particular any fixed rational linear combination with denominators
    coprime to $q_{\mathrm{rad}}$), or multiplying $y$ modulo $q_{\mathrm{rad}}^b$ by a fixed rational unit prime to
    $q_{\mathrm{rad}}$;
  \item exact fixed-width signed binary addition, subtraction, sign change, bit-order reversal within a component, and
    multiplication or division by a power of two when guard bits prevent overflow or inexact truncation.
\end{enumerate}
\end{fact}

Second, Fact~\ref{fact:oai-networks} states the finite algebraic certificates of the two fixed linear circuits at $h=100$ and $m=h^3=10^6$. The strict inequalities $s_b<W_b m$ and $s_c<W_c m$, which say that each network's recursive rank sum falls short of the trivial product of its role count and the dimension~$m$, are the sole source of the sub-$n\log n$ exponent saving.

\begin{fact}[Finite linear networks with a rank deficit~{\cite[Props.~7--8]{OAI26}}]\label{fact:oai-networks}
At ground-set parameter $h=100$ and dimension $m=h^3=10^6$, there exist two fixed finite linear circuits (over the
three-element subsets of $[h]$) with the following certified role counts, recursive rank sums, and endpoint actions for
arbitrary initial values on all auxiliary scratch roles:
\begin{enumerate}[label=(\alph*),nosep,leftmargin=*]
  \item \emph{Bit shear network}: a fixed pointwise XOR circuit on $W_b$ roles equipped with rational $m\times m$
    terminal and gate matrices whose completed scalar action is a role permutation $\rho_b$ satisfying
    $M_{\mathrm{out}(\rho_b(w))}-M_{\mathrm{in}(w)}=I_m$ for every input role $w$ (implementing the common address shear
    $(H,D)\mapsto(H+D,D)$ on $m$-tuples of address fields). Its role count $W_b$ and total edge-difference rank
    $s_b=\sum_e\operatorname{rank}_{\Q}(M_{\mathrm{head}(e)}-M_{\mathrm{tail}(e)})$ are
    \[
      W_b=177\,176\,569\,091\,445\,000\,000,\qquad
      s_b=177\,176\,569\,088\,785\,861\,287\,000\,000<W_b m,
    \]
    giving relative rank deficit $\eta_b=(W_b m-s_b)/(W_b m)=339/22\,587\,335\,000\,000$.
  \item \emph{Complex phase network}: a fixed scalar linear circuit on $W_c$ roles with coefficients in
    $\{0,\pm 1,\pm\tfrac12\}$ and nested binary subspace labels in $\mathbb F_2^m$ whose edge-residual orthonormal bases
    have total dimension $s_c$, implementing $C^{\otimes m}$ with $C=\frac{1+i}{2}I+\frac{1-i}{2}X$ on every role after
    prescribed source/sink diagonal phases and bank restoration, with
    \[
      W_c=1\,873\,807\,244\,643\,542\,670\,000,\qquad
      s_c=1\,873\,807\,244\,636\,671\,267\,308\,000\,000<W_c m,
    \]
    giving relative rank deficit $\eta_c=(W_c m-s_c)/(W_c m)=73/19\,906\,842\,167\,500$.
\end{enumerate}
\end{fact}

\paragraph{What the written proof depends on.}
No result of~\cite{HvdH21} is used anywhere in our proofs. Our written proof of Theorem~\ref{thm:main} imports only Facts~\ref{fact:oai-streams} and~\ref{fact:oai-networks}, corresponding to~\cite[Lemmas~3--4 and Props.~7--8]{OAI26}, as black boxes. Every other step is proved in Sections~\ref{sec:bridge}--\ref{sec:main} or re-derived in Appendix~\ref{app:parameter-transfer} by following~\cite[\S\S4--5]{OAI26} with all exponents and guard parameters exposed. Lemmas~15--21 and Proposition~18 of~\cite{OAI26}, covering synthetic ring transforms, chunk-swap layout, packed products, ring convolution, Gaussian resampling, and Bluestein's chirp, are unused on the main route and appear only for comparison in Appendix~\ref{app:bluestein}. The Lean~4 proof of Theorem~\ref{thm:lean} uses neither black box, constructing and verifying both $h=100$ networks, their algebraic certificates, and all Turing machine stream procedures from scratch.

\paragraph{Parameter-exposed tape procedures.}
As printed, \cite[Lemma~12 and Prop.~14]{OAI26} fix $\tau=\sigma=1-2^{-50}$, $\beta=\tfrac12$, and $C_1=20$. In Appendix~\ref{app:parameter-transfer} we follow~\cite[\S\S4--5]{OAI26} with all structural parameters exposed, deriving two procedures from Facts~\ref{fact:oai-streams} and~\ref{fact:oai-networks} together with Lemma~\ref{lem:transfer-long-record} for long-record address scans:
\begin{itemize}[nosep,leftmargin=*]
  \item \emph{Arbitrary-width interchange} (Lemma~\ref{lem:transfer12}): for every fixed rational $\tau\in(0,1)$ with $s_b/W_b<m^\tau$, two disjoint $u$-bit address fields are interchanged across arbitrary spectator fields in $O(V u^\tau)$ steps, and in $O(V(1+u^\tau))$ steps when their widths differ by one bit.

  \item \emph{Simultaneous butterfly layer} (Proposition~\ref{prop:prop14-param}, via packed selected-bit addition in Lemma~\ref{lem:transfer13}). Fix rational exponents $\tau,\sigma,\beta,c,\eps,\lambda,\lambda'$ and an integer guard exponent $C_1\ge\nu_{\mathrm{net}}+3=14$ with $\eps C_1<1$ satisfying the network bounds and admissibility constraints~\eqref{eq:walsh-layer-a}--\eqref{eq:walsh-layer-b} of Proposition~\ref{prop:prop14-param}, with $B_{\mathrm{net}},\nu_{\mathrm{net}},C_0$ as in~\eqref{eq:walsh-guard-constants}. Under the size bands~\eqref{eq:walsh-layer-bands} with $K=\lfloor d^c\rfloor$, the layer applies $H_0$ simultaneously to one common bit offset of $D\le d$ consecutive $K$-bit address chunks on a disk array of $\Rr_r$-records at grid precision~$q$. It evaluates the exact layer~$H$ inside $\Delta=\lceil C_0 d^{C_1}\rceil=o(q)$ internal guard bits, truncates once, and returns the disk array $Q_q(Hz)$ with $\|Q_q(Hz)-Hz\|_\infty<\sqrt2\cdot 2^{-q}$ in $O(V d^{\lambda'})$ steps.
\end{itemize}
Using the rational bound $\ln 10^6<14$ in place of $\ln 10^6<20$, Proposition~\ref{prop:deficits} certifies the sharper network exponents $\tau=1-2^{-40}$ and $\sigma=1-2^{-42}$ from the integers in Fact~\ref{fact:oai-networks}. Independently, the base-$10$ comparison $m^{10}<B_{\mathrm{net}}<m^{11}$ gives $\nu_{\mathrm{net}}=11$, certifying guard exponent $C_1=\nu_{\mathrm{net}}+3=14$ as noted in~\cite[\S5.4]{OAI26}.
